\scanpage{46}
\mspage{42}
% archive sheet number: 45

donc $T\cap T'=\varnothing$, donc $S\leq\varnothing$ donc $S=\varnothing$.

Cela impose donc une condition \inserted{nécessaire} en termes des \deleted{supports} \inserted{\unsure{déployements}}, pour qu'on ait
\[
X\ll\deleted{\unsure{X}}Y,
\]
c'est la condition suivante :
\[
\begin{aligned}
&\forall S\in\unsure{\mathrm{Dep}|\mathrm{D\acute{e}p}}\,X,\quad\exists T\in\mathrm{Dep}\,Y,\ \inserted{\text{(néc. unique, \unsure{not|tel} }T=\unsure{\psi|\tau}(S)\text{)}}\ \text{avec}\\
&S\leq T.
\end{aligned}
\tag*{\uline{Raff}}
\]
De plus, l'application
\[
\psi:\mathrm{Dep}(X)\mathbin{-}\mathrm{Dep}(Y)
\]
est croissante.

De plus, on a $X\unsure{\overset{\circ}{\ll}|\overset{\circ}{\Subset}}Y$ \unsure{sss|ssi} l'application $\psi$ précédente transforme le plus grand élément $X^{\circ}$ en le plus grand élément $Y^{\circ}$, i.e. si et seulement si on a
\[
\deleted{\unsure{X^{\circ}\ll}}\qquad X^{\circ}\leq Y^{\circ}
\]
(\uline{NB} \unsure{il est} entendu cependant qu'on sait déjà que $X\ll Y$\dots).

% margin note written sideways
% slanted
\margin{Si l'on a $X\trianglelefteq Y$, alors $\psi:\mathrm{Dep}\,X\longrightarrow\mathrm{Dep}\,Y$ est une inclusion d'ensembles ordonnés $\mathrm{Dep}(X)\subset\mathrm{Dep}(Y)$, et $\mathrm{Dep}(X)$ est une partie \unsure{ordonnée} \uline{fermée} de $\mathrm{Dep}\,Y$.}

\uline{Définition}. On dit qu'un \deleted{\unsure{magasin}} \inserted{atelier} est \uline{spatial}, si la condition nécessaire précédente Raff, sur $\mathrm{D\acute{e}p}\,X$ et $\mathrm{D\acute{e}p}\,Y$, pour que l'on ait $X\ll Y$,
